\subsection{Quadratic stars with center--leaf constraints}\label{sec:star}

Theorem~\ref{thm:polytope} is limited to small dimension, and
Section~\ref{sec:composition} shows that the blocks worth merging are
often stars: several pairs that share one variable. For stars there is an
exact algorithm whose cost is nearly linear in the size of the block.

Let $y$ be the center and $x_1,\dots,x_k$ the leaves, and consider
\begin{equation}\label{eq:star}
q(y,x)=a_0+a_1y+a_2y^2+\sum_{i=1}^k\bigl(d_ix_i^2+(e_iy+f_i)x_i\bigr)
\end{equation}
with rational coefficients of arbitrary signs. The domain $S$ is given by
finite rational bounds on all variables and by rational affine rows, each
of which involves $y$ and \emph{at most one} leaf; an equality is given as
two rows. Every scalarization $a\T(y,x)+b\T F(y,x)$ of a vector $F$ of
quadratics with this interaction pattern has the form~\eqref{eq:star}.
Rows or products that couple two leaves are excluded, and for good reason:
minimizing a concave separable quadratic over a box and a single row that
couples all variables is already NP-hard \citep{MoreVavasis1990}.

For fixed $y$, the rows of leaf $i$ restrict $x_i$ to an interval
$[L_i(y),U_i(y)]$, where $L_i$ is the maximum of the affine lower bounds and
$U_i$ the minimum of the affine upper bounds that these rows and the leaf
bounds define. Thus $L_i$ is convex and $U_i$ concave, both piecewise
affine. The set $J_i=\{y:L_i(y)\le U_i(y)\}$ is an interval because $L_i-U_i$
is convex. Given $y$, the leaves are constrained independently, so the
projection of $S$ onto $y$ is the interval $I$ obtained by intersecting
the bounds of $y$, the rows in $y$ alone, and all $J_i$. If $I$ is empty,
so is $S$.

\begin{theorem}[Exact support on constrained stars]\label{thm:star}
Let $m$ be the number of rows that involve a leaf, and suppose $S$ is
nonempty. Then $\min_Sq$ is attained at a rational point, and it can be
computed, together with a minimizer and a partition of $I$ into at most
$2m+4k+1$ intervals on each of which every leaf has an affine minimizer,
with $O((m+k)\log(m+k))$ arithmetic operations and comparisons. All
numbers involved have bit length polynomial in the input length; the
breakpoints have bit length bounded independently of $k$.
\end{theorem}

\begin{proof}
Define $\varphi_i(y)=\min\{d_ix_i^2+(e_iy+f_i)x_i:L_i(y)\le x_i\le U_i(y)\}$
for $y\in I$, so that $\min_Sq=\min_{y\in I}V(y)$ with
$V(y)=a_0+a_1y+a_2y^2+\sum_i\varphi_i(y)$. Let $s_i=m_i+2$ be the number of
affine functions that define $L_i$ and $U_i$, where $m_i$ is the number of
rows of leaf $i$; together $L_i$ and $U_i$ have at most $s_i-2$ breakpoints.

If $d_i>0$, the minimizer is $\xi_i=\max\{L_i,\min\{v_i,U_i\}\}$ with the
affine function $v_i(y)=-(e_iy+f_i)/(2d_i)$. The function $v_i-L_i$ is
concave and $v_i-U_i$ convex, so each has at most two isolated zeros, and
$\xi_i$ is affine between consecutive points of a set of at most $s_i+2$
breakpoints. If $d_i\le0$, a minimizer is an endpoint of the interval, and
\begin{equation}\label{eq:endpoint-factor}
\psi_i(U_i)-\psi_i(L_i)=(U_i-L_i)\bigl(d_i(U_i+L_i)+e_iy+f_i\bigr),
\qquad\psi_i(t)=d_it^2+(e_iy+f_i)t .
\end{equation}
The first factor is nonnegative on $I$. On each of the at most $s_i-1$
intervals where $L_i$ and $U_i$ are affine, the second factor is affine and
has at most one zero, which decides the preferred endpoint; this gives at
most $2s_i-3$ breakpoints. Taking all breakpoints together splits $I$ into
at most $\sum_i2s_i+1=2m+4k+1$ intervals on which every leaf rule is affine
and $V(y)=q(y,\xi(y))$ is a quadratic polynomial in $y$.

Each $\varphi_i$ is continuous on $I$, since it is the minimum of a
continuous function over an interval whose endpoints depend continuously
on $y$. Hence $V$ agrees with the quadratic of each interval on its
closure, and the minimum of $V$ is the least value of these quadratics at
the interval endpoints and, where the leading coefficient is positive, at
interior stationary points. All breakpoints are zeros of affine functions
with rational coefficients, so every candidate is rational, and the
minimizing $y$ together with $x_i=\xi_i(y)$ is a rational minimizer.

For the operation count, the upper envelope of $s_i$ lines is computed by
sorting slopes and one stack-based scan in $O(s_i\log s_i)$ operations, and
the breakpoints and rules of leaf $i$ in $O(s_i)$ more. Sorting all
breakpoints costs $O((m+k)\log(m+k))$. A sweep over the sorted breakpoints
maintains the three coefficients of $V$; at each breakpoint only the
leaves whose rule changes are updated, so the sweep takes $O(m+k)$
operations. Bit lengths are bounded in Appendix~\ref{app:star}.
\end{proof}

The counts in the proof are attained. A single concave leaf whose bounds
zigzag forces $2s_i-3$ rule changes, and $k$ hinge leaves force $k$
more, so the partition can have $\Theta(m+k)$ intervals. On instances
with $m=\Theta(k)$, the operation bound is optimal up to a constant factor
for algebraic computation trees: deciding whether the minimum is at least
a given value requires $\Omega((m+k)\log(m+k))$ operations, by Ben-Or's
theorem applied to a family whose decision set has $r!$ components
(Appendix~\ref{app:star}).

\paragraph{Relation to forest algorithms.}
\citet{DelPiaKhajavirad2026} solve box-constrained quadratic programs on
forests with $O(n^2)$ arithmetic operations. After an affine scaling of
each variable to $[0,1]$, a box-constrained star is such a forest, and
their dynamic program rooted at the center specializes to the
box-constrained case of Theorem~\ref{thm:star}; in that case our
contribution is only the $O(k\log k)$ count, which replaces their
$O(k^2)$ root merge. Rows that couple the center with a leaf are outside
their model. A shear $x_i\mapsto x_i-\delta y$ turns a leaf domain that is a
parallelogram with two vertical sides back into a box and keeps the star
structure, but triangles, trapezoids with nonparallel sides and polygons
with more sides cannot be reduced in this way. Such rows also break the
structural property on which their algorithm rests: the conditional value
of a leaf is then no longer a minimum of affine functions of $y$ over a
fixed set, and it can have convex kinks. For example, for $x\in[0,1]$
with rows $x\ge y-\tfrac12$ and $x\ge\tfrac12-y$, the minimum of $x$ is
$\abs{y-\tfrac12}$. The SOC descriptions of star hulls by
\citet{DeyKhajavirad2025} assume sign patterns of the square terms and
box domains, and the LP approximations of \citet{BienstockMunoz2018}
allow constraints but are approximate. We know of no earlier exact
algorithm for nonconvex stars with center--leaf rows; the construction
itself, conditioning on the center and solving parametric scalar
problems, is elementary.

\begin{remark}[Deeper trees]\label{rem:trees}
The argument uses that every leaf is a leaf. On the path
$x_1$--$x_2$--$x_3$ over $[0,1]^3$ with
$f=(x_2-\tfrac18)x_1+x_2^2-x_2+x_2x_3+x_3^2-x_3$, the conditional value of the
subtree $\{x_1,x_2\}$ given $x_3=t$ is
$\min\{-\tfrac18,\,-\tfrac14(1-t)^2\}$, which switches at the irrational
point $t=1-\sqrt2/2$, although $\min f=-\tfrac38$ is rational
(Appendix~\ref{app:star}). The rational partition above therefore does
not extend by nesting conditional values. \citeauthor{DelPiaKhajavirad2026}
handle such breakpoints for boxes with an explicit algebraic
representation; we do not know whether edge-local affine rows can be
added to deeper trees with a polynomial algorithm.
\end{remark}

The certificate of a star support value records the partition of $I$,
the affine rule of every leaf on every piece, the quadratic of $V$ on
every piece and the candidates; a checker rebuilds the partition from the
rows and coefficients and compares. Our implementation is simpler than
the algorithm of the proof: it intersects all pairs of envelope lines and
recomputes every leaf rule on every piece, which costs $O((m+k)^3)$
operations and stores $\Theta(k)$ rules per piece. This is adequate for the
blocks of Section~\ref{sec:computations}, where $k\le4$.
